\paragraph{IV 20.9: changing the projective immersion (cote 25, sites 55--57).}
Section 9 occupies sites 55--57. Paragraph 9.1 introduces the immersion $u_n:P\to P(n)$ associated with the degree-$n$ symmetric power and describes replacing $f$ by $u_nf$. Its examples refer back to 8.12.b and to singular hyperplane sections of a quadric; the latter passage has a handwritten replacement by ``dans $\mathbf P^3$''. The displayed direct-image expression has no visible equality before the symmetric power, and the edition preserves that gap. Proposition 9.2 is stated over $S=\operatorname{Spec}(k)$ for smooth $X$ and an unramified $f:X\to P$, with $n\geq2$. It records the conditions cited as equivalent in 8.8 and describes the singular points of the corresponding sections outside an exceptional locus of codimension at least two. The surrounding revision replaces the original single-point wording by a finite-number formulation and retains a separate assertion for an immersion. The note below it explicitly calls for a formulation of 8.8 that accommodates the unramified case, referring to condition (ii bis).

Proposition 9.3 and its proof run from site 56 into site 57. The opening has a typed insertion imposing geometric irreducibility of $X$ and a handwritten $n\geq2$; the dimension assumption concerns $f(X)$. Its three clauses concern sections through a given point, sections tangent there to a specified linear subvariety, and sections tangent to $X$. The geometric-irreducibility conclusions are proved in the draft by choosing polynomials $P,Q,R$ and examining $Q/P$ and $R/P$ in the function field. The proof writes out the degree-two choices and the exceptional planar discussion, followed by a marginal set of squared-coordinate choices for $\dim L\geq3$. That marginal note is linked to the body by matching oblique marks. Several overprinted coordinates and the handwritten point symbol remain doubtful in the transcription. Site 57 then proposes combining the argument with a reference beginning 8.18, whose final overprinted character is unresolved, under the unfinished heading Corollaire 9.4. A further appeal to 9.2 introduces the proposed ``cas excellent'' statement numbered 9.5, with a handwritten Th\'eor\`eme over a partly obscured typed title and a final note about a fundamental-group-like symbol and singular points. The edition preserves these drafting instructions and their uncertainty; it supplies no completion to either unfinished statement.
